%点群を分類するディープラーニングの関連研究

%1st paragraph
%1 分類の説明
Classification estimates which category a given data sample belongs to by associating it with the highest estimated probability.
%2 深層学習による点群分類の潮流
Deep learning models for point cloud classification have attracted interest \cite{grandini2020metrics}.

% %2nd paragraph
% %1
% While early models used voxelization \cite{Wu_2015_CVPR,7353481,Qi_2016_CVPR}, Qi et al. presented PointNet, which directly processes point clouds by considering permutation invariance \cite{grandini2020metrics}.
% %2
% PointCNN addresses convolution challenges on point clouds using the X-Conv operator to weight and permute input features \cite{NEURIPS2018_f5f8590c}.
% %3
% ShellNet extracts features from concentric shells to maintain local structures \cite{Zhang_2019_ICCV}, and A-CNN utilizes annular convolutions for hierarchical feature processing \cite{Komarichev_2019_CVPR}.
% %4
% SO-Net models spatial distributions via self-organizing maps (SOM) \cite{Li_2018_CVPR}, while Kd-network employs tree-based partitions for classification without grid reliance \cite{8237361}.

%2nd paragraph
%1PointNet以前の点群の分類
Early approaches relied on voxelization \cite{Wu_2015_CVPR,7353481,Qi_2016_CVPR}.
%2PointNet
Qi et al. presented PointNet \cite{grandini2020metrics}.
    %2-1深層学習モデルとしての位置づけ
    The network is a deep learning model for point cloud classification.
    %2-2共有MLPによる点ごとの特徴抽出
    PointNet applies shared multilayer perceptrons (MLPs) to each point in a set.
    %2-3max poolingによる特徴集約
    The model aggregates per-point features using max pooling.
    %2-4置換不変性
    The architecture enforces permutation invariance over the input order.

 \begin{deletedblock}
 %3rd paragraph
%1CNNを使った点群の分類
Point cloud classifiers often adopt convolutional neural networks (CNNs).
  %1-1PointCNN
  Li et al. presented PointCNN \cite{NEURIPS2018_f5f8590c}.
    %1-1-1X-Convの導入
    The method introduces X-Conv for point cloud deep learning.
    %1-1-2置換不変畳み込み
    The operator performs end-to-end permutation-invariant convolution.
    %1-1-3座標からの階層的局所特徴学習
    The architecture learns hierarchical local features from point coordinates.
    %1-1-4点属性の利用
    The architecture uses point attributes as additional inputs.
  %1-2ShellNet
  Zhang et al. presented ShellNet \cite{Zhang_2019_ICCV}.
    %1-2-1ShellConvによる近傍特徴集約
    The model aggregates neighborhood features using ShellConv.
    %1-2-2同心円シェル分割
    The operator partitions local point sets into concentric shells.
    %1-2-3局所構造の保持
    The network extracts hierarchical features to preserve local structures.
  %1-3A-CNN
  Komarichev et al. presented A-CNN \cite{Komarichev_2019_CVPR}.
    %1-3-1環状畳み込み
    The method applies annular convolutions on a three-dimensional (3D) point cloud.
    %1-3-2最遠点サンプリング
    The pipeline uses farthest point sampling (FPS) to select center points.
    %1-3-3法線入力
    The model takes point normals as inputs.
    %1-3-4座標入力
    The model takes point coordinates as inputs.
    %1-3-5近傍からの階層的特徴学習
    The network learns hierarchical features from sampled neighborhoods.

%4th paragraph
%1点群の分類におけるその他のアプローチ
  %1-1SO-Net
  Li et al. proposed SO-Net \cite{Li_2018_CVPR}.
    %1-1-1固定長符号化
    The method encodes a point cloud into a fixed-length representation.
    %1-1-2SOMによる空間分布モデル化
    The encoder builds a self-organizing map (SOM) over point coordinates \cite{58325}.
    %1-1-3k-NNによる特徴洗練
    The pipeline refines features using k-nearest neighbor (k-NN) search.
    %1-1-4max pooling
    The network applies max pooling over the refined features.
  %1-2Kd-Network
  Klokov et al. proposed Kd-network \cite{8237361}.
    %1-2-1kd-tree構築
    The architecture constructs a kd-tree over point coordinates.
    %1-2-2ノード間ベクトルへのアフィン変換
    The model applies affine transforms on inter-node vectors along the tree.
    %1-2-3MLPによる分類
    The pipeline passes the resulting features through an MLP for classification.
    %1-2-4グリッド非依存の分類
    The method classifies objects without rasterizing points onto a uniform grid.

% \end{spverbatim}
 \end{deletedblock}
%5th paragraph
%1 既存研究の特徴
The existing research shown above can be broadly classified into the following two categories:
\begin{enumerate}
  %1-1直接分類アーキテクチャ
  \item Architectures for direct point cloud classification (e.g., PointNet, PointCNN).
  %1-2局所関係の空間モデリング
  \item Spatial modeling techniques using local relationships (e.g., ShellNet, SO-Net, A-CNN, Kd-network).
\end{enumerate}
%2 本研究の特徴
Our research targets the fine-grained classification of hand movements in a multi-person office environment.
%3 本研究との差分 
Our research differs from the macro-posture or static object classification focused on in the categories above.
  % %3-1 本研究との差分
  % The study additionally proposes an edge system for capturing detailed hand movements using a LiDAR sensor network.
